EXERCISE (KERNELS COMMUTE WITH LIMITS). Suppose $\mathcal{C}$ is an abelian category, and $a:I\to\mathcal{C}$ and $b:I\to\mathcal{C}$ are two diagrams in $\mathcal{C}$ indexed by $I$. For convenience, let $A_i=a(i)$ and $B_i=b(i)$ be the objects in those two diagrams. Let $h_i:A_i\to B_i$ be maps commuting with the maps in the diagrams. (Translation: $h$ is a natural transformation of functors $a\to b$; see §1.1.21.) Then the $\ker h_i$ form another diagram in $\mathcal{C}$ indexed by $I$. Describe a canonical isomorphism
$$\varprojlim_i\ker h_i\cong\ker\left(\varprojlim_i A_i\to\varprojlim_i B_i\right),$$
assuming the limits exist. Implicit in the previous exercise is the idea that limits should somehow be understood as functors.
